\subsection{固定楕円境界の重み付き係数と判別式}
\label{subsec:cubic-CY-three-alpha-elliptic-fixed-boundary}

\begin{theorem}\label{thm:cubic-CY-three-alpha-elliptic-fixed-boundary-coefficients}
前節の滑らかなHesse先頭環 $S$ とその中央三次 $g$ を一つ固定し、
$C=S/(g)$ とする。associated gradedを $C$ と同定したfiltered deformationの
関手は、非被約底も含めてアフィン七次元空間で表される。
係数は、対角零の対称行列 $A$ の三成分、三定数 $b_i$ と中央切断値 $\lambda$ である。
先頭の同定を共通倍率の範囲で変える操作はそれぞれ重み $1,2,3$ で作用する。
従って零変形を除いたこの倍率商は重み付き射影stack
$\mathbb P(1,1,1,2,2,2,3)$ である。
\end{theorem}
\begin{proof}
$C_1$ の標準三生成元を持ち上げると、それらはfiltered生成元となる。
三二次関係のずれは線型式と定数である。
Rees環上で三次の巡回tensor恒等式を用い、中央次数一の $t$ で割ると、
線型係数の交代部分が二次交換子の線型関係を与える。
先頭の純平方係数が全て非零なので $S$ に純交換子関係はなく、
対応する三つの交換子は底上も独立である。従って $A$ は対称である。

持ち上げを $X_i=Y_i+u_i$ と変えると、線型行列は
$A-\operatorname{Jac}(f)(u)$ に変わる。ここで
$f_i(u)=(p+q)u_ju_\ell+r u_i^2$ である。
対角を零にする一意の選択は $u_i=A_{ii}/(2r)$ である。
従って指定した先頭同型に対する正規化は底上でも一意で、
持ち上げの選択に依存せず貼り合わさる。

この二次包絡は、全低次係数の被約普遍空間でのPBW定理と
次数別階数一定性により全次数で平坦である。
従って任意の底への引き戻しも平坦である。
Rees環への全射の三次核は、$t=0$ への制限で $g$ に同型に移る。
その一意の持ち上げ $G$ について、左右乗法は各次数で分裂単射であり、
Hilbert階数の比較から左右とも全核を生成する。
$G$ の正常化作用は、$t$ を固定して $S$ 上恒等になる。
従って $X_i\mapsto X_i+v_it$ の形である。
これが二次関係を保つには $2r v_i=0$ が必要なので恒等、
よって $G$ は中央である。

普遍空間における中央三次の持ち上げの全体は、
$t^3$ の直線による $g$ の直線の拡大である。
次数付きスケーリングに等変な持ち上げを選べる。
実際、普遍係数環は重み付き多項式環であり、
中央核の先頭 $g$ への全射に任意の持ち上げを取り、指定重み成分に射影すればよい。
この持ち上げを $\omega$ とすれば全切断は $\omega-\lambda$ となる。
一次生成元の倍率変更 $X_i\mapsto cX_i$ は、
係数を同時に $(A,b,\lambda)\mapsto(cA,c^2b,c^3\lambda)$ に移す規約で与えられる。
逆にこの七係数は平坦なRees商を作る。
任意の先頭恒等付き同型は生成元の平行移動であり、
対角零の正規化を保つなら恒等なので、全矢も確定する。
\end{proof}

\begin{proposition}\label{prop:cubic-CY-three-alpha-elliptic-discriminant-weight}
$\sigma$ の位数を無限とする。
固定先頭上の全低次係数にわたり、定理
\ref{thm:cubic-CY-three-alpha-elliptic-diagonal-test}の八次判別式は
一つの重み24の多項式となる。
\end{proposition}
\begin{proof}
これは対角零にする前の係数空間でまず証明してよい。
有限局所自由コチェインの重み $-3$ 以下の四次コホモロジーは
全点で八次元である。隣接微分の階数の和は一定で、
各階数の下界は開条件なので、両階数は局所一定である。
従って四次コホモロジーはrank8のベクトル束をなし、
その形成は底変更に両立する。
重み0まで上げた同じコチェインの商により中央元の作用も相対的に定まる。
従って特性多項式の係数は低次係数の多項式になる。
倍率変更で作用は三乗倍されるので、
根 $\lambda$ の重みを3とすればモニック八次式の全体は重み24である。
平行移動で正規化した七係数にも同じ結論が従う。
\end{proof}

\begin{theorem}\label{thm:cubic-CY-three-alpha-elliptic-eight-monodromy}
固定先頭 $(p,q,r)=(1,2,3)$ に対し、八次判別式の根を交換する
幾何学的モノドロミー群は $S_8$ である。
従って固定楕円境界の七係数空間上で、双加群非滑らか切断の重み24の因子は
幾何学的に既約である。
\end{theorem}
\begin{proof}
まず平行移動を規格化する前の対称係数と定数係数を全て動かす。
定理\ref{thm:cubic-CY-three-alpha-elliptic-fixed-boundary-coefficients}の
一意な平行移動で、これは規格化後の空間に三つの自由座標を加えただけであり、
切断値が多項式で平行移動することは根の群を変えない。

$A=\operatorname{diag}(1,2,4)$、$b=0$ で得た八次式は $\mathbb Q$ 上既約である。
素数5と17での分離な因数の次数はそれぞれ $(1,7)$ と $(2,3,3)$。
従って特殊化のGalois群は推移的で、7-cycleと互換を含む。
8点上の非自明なblockの大きさは2または4なので、7-cycleによりprimitiveである。
primitive群の互換の共役全体は、点を頂点、互換を辺とする連結グラフを与える。
従ってこれらは全対称群を生成する。
根の分離開上の根番号付けは有限エタールなので、この有理点のGalois群は
一般の算術群の部分群となる。従って算術群は $S_8$ である。

幾何学的群を決定するため、$A=\operatorname{diag}(1,2,a)$、$b=0$ の
直線上で、八次式の $\lambda$ に関する通常の判別式を取る。
重み24なので、$\lambda^j$ の係数の $a$ 次数は $3(8-j)$ 以下である。
従って整数 $a=0,\ldots,24$ での有理数係数と37での還元で、
モニック八次式のこの直線上での還元を一意に復元できる。
この25点で係数の分母は37と互いに素で、補間行列も37上可逆である。
その還元の $\lambda$ 判別式では
\[
 a^4+2a^3+22a^2+9a+21
\]
が既約で、重複度は3である。
ゆえに標数零の直線上でも判別式は $\overline{\mathbb Q}(a)$ の平方ではない。
実際、幾何学的な平方なら、$\mathbb Q[a]$ の全既約因子の重複度が偶数であり、
37で非零の還元も単元倍の平方となる。
一般の係数空間の判別式が平方でも、非零となるこの直線への制限は平方になる。

幾何学的群は算術群の正規部分群である。
判別式が平方ではないため奇置換を含み、従って $S_8$ 全体である。
特に根の被覆は幾何学的に連結で、モニック多項式の零点超曲面は幾何学的に既約である。
計算は
\texttt{verify\_cubic\_CY\_\allowbreak{}three\_alpha\_elliptic\_\allowbreak{}monodromy.py}
で再現できる。
\end{proof}
